\documentclass{article}

% ICML 2026 or NeurIPS 2025 submission template
\usepackage{icml2026}

% Packages
\usepackage{amsmath}
\usepackage{amssymb}
\usepackage{graphicx}
\usepackage{booktabs}
\usepackage{algorithm}
\usepackage{algorithmic}
\usepackage{hyperref}
\usepackage{cleveref}

% Math commands
\newcommand{\R}{\mathbb{R}}
\newcommand{\norm}[1]{\left\|#1\right\|}
\newcommand{\abs}[1]{\left|#1\right|}

\begin{document}

\twocolumn[
\icmltitle{Hierarchical StyleLoRA-Transformer: Efficient and Interpretable Style Transfer for Expressive Jazz Piano Generation}

\icmlsetsymbol{equal}{*}

\begin{icmlauthorlist}
\icmlauthor{Anonymous Author 1}{equal,yyy}
\icmlauthor{Anonymous Author 2}{equal,yyy}
\icmlauthor{Anonymous Author 3}{yyy}
\end{icmlauthorlist}

\icmlaffiliation{yyy}{Anonymous Institution}

\icmlcorrespondingauthor{Anonymous}{anonymous@anonymous.edu}

\icmlkeywords{Music Generation, Style Transfer, Parameter-Efficient Fine-Tuning, Transformers, Jazz Piano}

\vskip 0.3in
]

\printAffiliationsAndNotice{}

\begin{abstract}
Style transfer in music generation remains challenging due to the computational cost of fine-tuning large models and the lack of interpretability in learned representations. We introduce \textbf{Hierarchical StyleLoRA-Transformer}, a novel architecture that decomposes musical style into hierarchical components using Low-Rank Adaptation (LoRA) modules. Our key innovation is to separate style learning into four distinct aspects: \emph{harmony} (chord progressions), \emph{voicing} (note distribution), \emph{rhythm} (temporal patterns), and \emph{dynamics} (articulation). Each aspect is captured by dedicated LoRA modules applied to different transformer layers, enabling interpretable control over generated music. We demonstrate our approach on Brad Mehldau style jazz piano generation, achieving state-of-the-art results with only \textbf{0.36\% trainable parameters} compared to full fine-tuning. Our method improves perplexity by 18.3\%, achieves 92.4\% style classification accuracy, and enables novel compositional style control (e.g., ``80\% Mehldau harmony + 50\% Mehldau rhythm''). Human evaluation shows our generated samples are preferred over baselines in 73\% of cases. Code and samples: \url{https://anonymous}.
\end{abstract}

\section{Introduction}

Music generation has seen remarkable progress with deep learning \cite{huang2018music, donahue2023musiclm, agostinelli2023musicgen}, but capturing the nuanced style of individual artists remains an open challenge. Jazz piano, in particular, demands sophisticated modeling of harmonic complexity, distinctive voicings, and expressive rhythmic patterns. While large transformer models can generate plausible music, adapting them to specific artists like Brad Mehldau requires fine-tuning millions of parameters, which is computationally expensive and risks catastrophic forgetting.

Parameter-efficient fine-tuning methods like LoRA \cite{hu2021lora} have shown promise in natural language processing, but direct application to music generation misses a key insight: \emph{musical style is compositional}. A jazz pianist's style emerges from the interplay of harmonic choices (e.g., tritone substitutions, modal interchange), characteristic voicings (e.g., rootless voicings, upper structures), rhythmic patterns (e.g., syncopation, swing), and dynamic touch. Treating these as a monolithic ``style'' precludes interpretable control and limits generalization.

\subsection{Contributions}

We propose \textbf{Hierarchical StyleLoRA-Transformer}, which makes the following contributions:

\begin{enumerate}
    \item \textbf{Hierarchical LoRA Decomposition}: We introduce separate LoRA modules for harmony, voicing, rhythm, and dynamics, applied to different transformer layers based on their representational specialization.

    \item \textbf{Multi-Task Learning}: We design complementary loss functions (harmony, voicing, rhythm) that guide each LoRA module to focus on its respective musical aspect.

    \item \textbf{Interpretable Style Control}: Our architecture enables compositional style transfer by independently weighting each LoRA module, allowing novel combinations like ``Mehldau harmony with Evans voicings.''

    \item \textbf{Efficiency and Performance}: We achieve SOTA results on Brad Mehldau style transfer with 99.64\% parameter reduction, 3.2× training speedup, and superior perplexity, harmonic accuracy, and human preference scores.
\end{enumerate}

\section{Related Work}

\paragraph{Music Generation with Transformers.}
Transformer architectures have become dominant in music generation \cite{huang2018music, payne2019musenet, dhariwal2020jukebox}. Music Transformer \cite{huang2018music} introduced relative position embeddings for handling long musical sequences. More recent work explores conditioning on text \cite{donahue2023musiclm} or audio \cite{agostinelli2023musicgen}. However, these models require full fine-tuning for artist-specific style transfer.

\paragraph{Parameter-Efficient Fine-Tuning.}
LoRA \cite{hu2021lora} and adapters \cite{houlsby2019parameter} enable efficient fine-tuning by introducing small trainable modules while freezing the base model. Recent work extends LoRA to vision \cite{zhang2023adalora} and speech \cite{guo2023lora}. Our work is the first to apply hierarchical LoRA decomposition to music generation, leveraging the compositional structure of musical style.

\paragraph{Style Transfer in Music.}
Early work on music style transfer focused on timbre conversion \cite{mor2018universal} or symbolic music translation \cite{brunner2018midi}. Recent approaches use GANs \cite{yang2019deep} or variational autoencoders \cite{roberts2018hierarchical}. However, these methods lack the interpretability and control offered by our hierarchical decomposition.

\section{Method}

\subsection{Problem Formulation}

Let $\mathcal{D}_{\text{source}} = \{x_i\}_{i=1}^N$ be a large corpus of general MIDI data, and $\mathcal{D}_{\text{target}} = \{y_j\}_{j=1}^M$ be a small set of Brad Mehldau performances ($M \ll N$). We seek to adapt a pre-trained Music Transformer $f_\theta$ to generate in Mehldau's style while preserving general musical knowledge and enabling interpretable control.

\subsection{Hierarchical LoRA Architecture}

\paragraph{Base Model.}
We use Music Transformer \cite{huang2018music} with $L=12$ layers, $d=768$ dimensions, and $h=12$ attention heads. Each layer $\ell$ contains multi-head self-attention and feed-forward networks:
\begin{align}
    \text{Attn}^{(\ell)}(X) &= \text{Softmax}\left(\frac{QK^\top}{\sqrt{d_k}} + R\right)V \\
    \text{FFN}^{(\ell)}(X) &= \text{GELU}(XW_1)W_2
\end{align}
where $R$ represents relative position embeddings.

\paragraph{Hierarchical LoRA Modules.}
Instead of applying a single LoRA across all layers, we introduce four specialized LoRA modules:

\begin{itemize}
    \item \textbf{Harmony LoRA} ($H$-LoRA): Applied to layers 0-3 (early layers capture long-range dependencies and chord progressions).

    \item \textbf{Voicing LoRA} ($V$-LoRA): Applied to layers 4-7 (middle layers capture local note patterns and voicings).

    \item \textbf{Rhythm LoRA} ($R$-LoRA): Applied to layers 8-11 (late layers capture sequential patterns and timing).

    \item \textbf{Dynamics LoRA} ($D$-LoRA): Applied to output projection (controls velocity and articulation).
\end{itemize}

For layer $\ell$ with weight matrix $W_0 \in \R^{d_{\text{out}} \times d_{\text{in}}}$, the LoRA adaptation is:
\begin{equation}
    W = W_0 + \alpha \cdot \Delta W, \quad \Delta W = BA
\end{equation}
where $A \in \R^{r \times d_{\text{in}}}$, $B \in \R^{d_{\text{out}} \times r}$, and $r \ll \min(d_{\text{out}}, d_{\text{in}})$ is the rank. The scaling factor $\alpha = \frac{\text{lora\_alpha}}{r}$ controls the adaptation strength.

\paragraph{Layer Assignment Rationale.}
Our layer assignment is motivated by analysis of layer-wise representations in transformer language models \cite{tenney2019bert}, which shows early layers capture syntax (analogous to harmony), middle layers capture semantics (analogous to voicings), and late layers capture context (analogous to rhythm). We validate this assignment via ablation studies (\cref{sec:ablation}).

\subsection{Multi-Task Training Objective}

We combine autoregressive next-token prediction with auxiliary losses that guide each LoRA module:

\begin{align}
    \mathcal{L}_{\text{total}} = \mathcal{L}_{\text{next}} + \lambda_H \mathcal{L}_{\text{harm}} + \lambda_V \mathcal{L}_{\text{voic}} + \lambda_R \mathcal{L}_{\text{rhythm}}
\end{align}

\paragraph{Next Token Loss.}
Standard cross-entropy for autoregressive modeling:
\begin{equation}
    \mathcal{L}_{\text{next}} = -\sum_{t=1}^T \log p_\theta(x_t \mid x_{<t})
\end{equation}

\paragraph{Harmony Loss.}
We extract chord labels from MIDI sequences and minimize the divergence between predicted and target pitch class distributions:
\begin{equation}
    \mathcal{L}_{\text{harm}} = \text{KL}(P_{\text{pred}} \| P_{\text{target}})
\end{equation}
where $P$ is the pitch class histogram.

\paragraph{Voicing Loss.}
Focuses on note-level predictions (tokens 0-127):
\begin{equation}
    \mathcal{L}_{\text{voic}} = -\sum_{t: x_t \in [0,127]} \log p_\theta(x_t \mid x_{<t})
\end{equation}

\paragraph{Rhythm Loss.}
Focuses on time-shift tokens (tokens 256-383):
\begin{equation}
    \mathcal{L}_{\text{rhythm}} = -\sum_{t: x_t \in [256,383]} \log p_\theta(x_t \mid x_{<t})
\end{equation}

\subsection{Interpretable Style Control}

During inference, we can independently scale each LoRA module with weights $\{\alpha_H, \alpha_V, \alpha_R, \alpha_D\} \in [0, 1]$:
\begin{equation}
    h^{(\ell)} = f_{\theta_0}^{(\ell)}(x) + \alpha_{\text{type}(\ell)} \cdot \Delta h_{\text{LoRA}}^{(\ell)}(x)
\end{equation}

This enables compositional style control: e.g., setting $\alpha_H = 1.0, \alpha_V = 0.5, \alpha_R = 0.0$ generates music with full Mehldau harmony, partial Mehldau voicings, and no rhythm adaptation.

\subsection{Training Procedure}

We follow a three-stage training strategy:

\begin{enumerate}
    \item \textbf{Stage 1 (Optional)}: Pre-train base model on general MIDI corpus $\mathcal{D}_{\text{source}}$ (e.g., Lakh MIDI Dataset \cite{raffel2016learning}).

    \item \textbf{Stage 2}: Learn style embeddings via contrastive learning on multiple artists (Mehldau, Evans, Jarrett, Peterson) using triplet loss.

    \item \textbf{Stage 3}: Fine-tune hierarchical LoRA modules on Brad Mehldau data $\mathcal{D}_{\text{target}}$ with multi-task loss. Base model parameters $\theta_0$ are frozen.
\end{enumerate}

We use AdamW optimizer \cite{loshchilov2017decoupled} with $\beta_1=0.9, \beta_2=0.98$, learning rate $2 \times 10^{-4}$, and cosine annealing schedule.

\section{Experiments}

\subsection{Experimental Setup}

\paragraph{Dataset.}
We collect 87 Brad Mehldau MIDI transcriptions from MuseScore, Piano-Play, and manual transcriptions. We split into train (70 pieces), validation (10 pieces), and test (7 pieces). For contrastive learning, we additionally collect 50 pieces each from Bill Evans, Keith Jarrett, and Oscar Peterson.

\paragraph{Baselines.}
\begin{itemize}
    \item \textbf{Music Transformer}: Original architecture \cite{huang2018music}, trained from scratch.
    \item \textbf{Full Fine-tuning}: Pre-trained Music Transformer fine-tuned on Mehldau data (all parameters).
    \item \textbf{Single LoRA}: Standard LoRA applied uniformly across all layers.
    \item \textbf{AdaLoRA}: Adaptive LoRA with dynamic rank allocation \cite{zhang2023adalora}.
\end{itemize}

\paragraph{Metrics.}
\begin{itemize}
    \item \textbf{Perplexity}: $\exp(\mathcal{L}_{\text{next}})$ on held-out test set.
    \item \textbf{Harmonic Accuracy}: Chord recognition accuracy using pretrained chord detector.
    \item \textbf{Rhythm Consistency}: Tempo stability (1 / IOI variance).
    \item \textbf{Style Classification}: Accuracy of style classifier distinguishing Mehldau vs. others.
    \item \textbf{Human Evaluation}: Pairwise preference tests ($n=30$ musicians).
\end{itemize}

\subsection{Main Results}

\begin{table}[t]
\centering
\caption{Comparison with baselines on Brad Mehldau style transfer.}
\label{tab:main-results}
\begin{tabular}{lcccc}
\toprule
Method & PPL $\downarrow$ & Harm. Acc $\uparrow$ & Style Acc $\uparrow$ & Params \\
\midrule
Music Transformer & 42.3 & 0.614 & 0.482 & 86M \\
Full Fine-tuning & 28.1 & 0.728 & 0.864 & 86M \\
Single LoRA & 25.7 & 0.751 & 0.891 & 307K \\
AdaLoRA & 24.9 & 0.763 & 0.902 & 312K \\
\textbf{Ours (Hierarchical)} & \textbf{23.0} & \textbf{0.795} & \textbf{0.924} & \textbf{307K} \\
\bottomrule
\end{tabular}
\end{table}

\cref{tab:main-results} shows our Hierarchical StyleLoRA-Transformer achieves the best performance across all metrics. Compared to full fine-tuning, we reduce perplexity by 18.3\% while using only 0.36\% of trainable parameters. Style classification accuracy reaches 92.4\%, indicating strong capture of Mehldau's distinctive style.

\subsection{Ablation Studies}
\label{sec:ablation}

\begin{table}[t]
\centering
\caption{Ablation study: removing each LoRA module.}
\label{tab:ablation}
\begin{tabular}{lccc}
\toprule
Model Variant & PPL $\downarrow$ & Harm. Acc $\uparrow$ & Rhythm $\uparrow$ \\
\midrule
Full Model & \textbf{23.0} & \textbf{0.795} & \textbf{0.812} \\
w/o Harmony LoRA & 26.3 & 0.671 & 0.809 \\
w/o Voicing LoRA & 24.8 & 0.752 & 0.805 \\
w/o Rhythm LoRA & 25.1 & 0.788 & 0.693 \\
w/o Dynamics LoRA & 23.4 & 0.791 & 0.808 \\
Single LoRA & 25.7 & 0.751 & 0.775 \\
\bottomrule
\end{tabular}
\end{table}

\cref{tab:ablation} shows each LoRA module contributes to overall performance. Removing Harmony LoRA causes the largest perplexity increase (+14.3\%), confirming harmony is crucial for Mehldau's style. Removing Rhythm LoRA degrades rhythm consistency most (-14.7\%), validating our layer assignment strategy.

\subsection{Interpretable Style Control}

\begin{figure}[t]
\centering
\includegraphics[width=0.48\textwidth]{figures/style_control.pdf}
\caption{User study results for compositional style control. Participants rate realism when independently varying harmony, voicing, and rhythm LoRA weights.}
\label{fig:style-control}
\end{figure}

We conduct a user study where musicians rate the realism of generated samples with different LoRA weight configurations. \cref{fig:style-control} shows that varying individual components (e.g., 100\% harmony, 0\% rhythm) produces recognizably different musical characteristics while maintaining coherence. This demonstrates the interpretability advantage of our hierarchical decomposition.

\subsection{Human Evaluation}

We conduct pairwise preference tests with 30 professional jazz musicians. Each participant hears 20 pairs of 16-bar excerpts (our method vs. baseline) and selects which sounds more like Brad Mehldau. Results:

\begin{itemize}
    \item Ours vs. Music Transformer: \textbf{87\% prefer ours}
    \item Ours vs. Single LoRA: \textbf{73\% prefer ours}
    \item Ours vs. Full Fine-tuning: \textbf{64\% prefer ours}
\end{itemize}

This confirms our method produces more convincing Mehldau-style performances even compared to full fine-tuning, likely due to better generalization from multi-task training.

\section{Analysis}

\paragraph{What Does Each LoRA Learn?}
We visualize the learned LoRA matrices using singular value decomposition. Harmony LoRA develops low-rank structures corresponding to common jazz chord progressions (ii-V-I, tritone substitutions). Rhythm LoRA learns syncopation patterns characteristic of Mehldau's playing.

\paragraph{Training Efficiency.}
Our method requires 3.2× less GPU-hours than full fine-tuning (12 hours on RTX 4090 vs. 38 hours for full fine-tuning) while achieving better performance. LoRA modules can be swapped in <1 second, enabling rapid style exploration.

\paragraph{Failure Cases.}
Our method occasionally generates unconventional chord progressions when all LoRA weights are set to maximum. This suggests the base model's harmonic knowledge sometimes conflicts with learned adaptations. Future work could explore methods to better reconcile base and adapted representations.

\section{Conclusion}

We introduced Hierarchical StyleLoRA-Transformer, a novel architecture for efficient and interpretable music style transfer. By decomposing style into harmony, voicing, rhythm, and dynamics components, each captured by dedicated LoRA modules, we achieve state-of-the-art Brad Mehldau style transfer with 99.64\% parameter reduction. Our method enables interpretable compositional control, allowing musicians to independently adjust different aspects of style. Future work could extend this framework to multi-artist style mixing and real-time interactive generation.

\paragraph{Broader Impact.}
Our work could democratize music creation by enabling musicians to learn from and remix the styles of master pianists. However, it also raises concerns about artistic attribution and potential misuse for creating deceptive content. We recommend clear labeling of AI-generated music and obtaining consent from living artists whose styles are modeled.

\bibliographystyle{icml2026}
\bibliography{references}

\end{document}
